\Needspace{18\baselineskip}
\section{Independent local incidence semantics for search}
\label{sec:local-incidence}
This section proves that erasure-based state search coincides with a presentation
that inspects adjacency in the original graph and specifies the target by explicit
incidence equations. Neither the local interaction predicate nor the rewiring
structure mentions erasure or coordinate search. The result applies to arbitrary
states, including states with cycles, and preserves ordered arguments and complete
non-focus labels.

The incidence correspondence alone does not establish contextual-rewrite
adequacy. Section~\ref{sec:contextual-search} now completes the additional
construction of finite interaction patterns and proves their factorization
and plugging correspondence through explicit focus-free simple hole contexts.

\subsection{Interactions in the original graph}
\begin{definition}[Local interaction]\label{term:local-interaction}
Let $f$ be a focus of a graph $C$ with no output boundary. A local interaction
selects a pair $(m,w)$ consisting of the next mode and an original graph port.
The variable rule requires search mode, a variable-labelled edge, and that
edge's source list to be exactly $[f.\mathsf{target}]$; it selects rewrite mode
at $f.\mathsf{source}$. Entry requires search mode and an operation whose
source equals $f.\mathsf{target}$, and uses its ordered entry selection.
Return requires backtracking mode and an argument index $i$ whose lookup is
$f.\mathsf{source}$, and uses the operation's ordered return selection.
Entry and return include both the next-argument and finished-operation cases.
\end{definition}
The selected port is not required to differ from the old focus target. For
example a nullary operation immediately adjacent to the focus selects its own
source, which is that target. The joining operation below handles this case.
\lean{Graph/LocalInteraction}{10}{19}{Linear.LocalInteraction}

\begin{lemma}[Erasure and operation views]
An operation edge differs from the focus edge. Its erasure retains its label,
contracts its source and all argument and thunk ports, and is again an operation
view. Erasure commutes with entry and return selection. Conversely every erased
operation has an original operation view with the same edge and operation label,
and with precisely those contracted source and argument data.
\end{lemma}
\begin{proof}
Operation and focus labels are disjoint at every depth, proving the edge
inequality. Mapping contraction over the singleton source and concatenated target
lists proves the erased view's incidence equations; list length is unchanged.
For entry, split on whether the argument list is empty. For return, commute
lookup with list mapping and split on the next-argument lookup. Finally, recover
an original operation view from the surviving edge's label. Uniqueness of the
source and ordered argument data for a given edge identifies its erasure with
the supplied view.
\end{proof}
\lean{Graph/LocalInteraction}{24}{27}{Linear.SearchOperation.ne_focus}
\lean{Graph/LocalInteraction}{30}{39}{Linear.SearchOperation.erase}
\lean{Graph/LocalInteraction}{41}{43}{Linear.SearchOperation.erase_enter}
\lean{Graph/LocalInteraction}{45}{49}{Linear.SearchOperation.erase_resume}
\lean{Graph/LocalInteraction}{51}{59}{Linear.SearchOperation.unerase}

\begin{theorem}[Adjacency correspondence]
A coordinate step from the erased focus coordinate to $q$ exists exactly when
an original-graph local interaction selects $(m,w)$ and
$q=(m,\mathsf{contract}(w))$.
\end{theorem}
\begin{proof}
For the forward construction from local interactions, erase the interacting
edge. A variable or operation source adjacent to the focus target contracts to
the erased focus position. A returning argument equal to the focus source
contracts to that same position. Apply the corresponding coordinate rule and
the operation-selection commutation equations.

For the converse, analyze the coordinate rule. In the variable case,
erased-source membership reflects to adjacency at the original focus target;
the variable signature makes the original source list a singleton. In the entry
case, recover the original operation view and use the same source-membership
reflection. In the return case, invert lookup in the mapped argument list to
recover the original argument at the same index. A non-focus target incidence
cannot equal the deleted focus target. Thus lifting its contraction recovers
the original argument, and the coordinate equality forces this argument to be
the focus source. Uniqueness of operation selections gives the required target
coordinate in both operation cases.
\end{proof}
\lean{Graph/LocalInteraction}{66}{96}{Linear.LocalInteraction.sound}
\lean{Graph/LocalInteraction}{100}{147}{Linear.LocalInteraction.complete}
\lean{Graph/LocalInteraction}{150}{157}{Linear.LocalInteraction.iff_search}

\subsection{Independent equations for the target graph}
\begin{definition}[Joining and moving ports]\label{term:local-port-movement}
Write $a$ and $b$ for the old focus source and target. Define functions on the
original graph's ports by
\[
 J(v)=\begin{cases}a&v=b,\\v&v\ne b,\end{cases}
 \qquad
 M_w(v)=\begin{cases}b&J(v)=J(w),\\J(v)&J(v)\ne J(w).\end{cases}
\]
The first joins the old focus endpoints. The second redirects source incidences
at the selected joined port to the reusable focus target. These equations use
only original ports and equality, without constructing an erased graph.
\end{definition}
\lean{Graph/LocalRewiring}{12}{14}{Linear.Focussed.joinPort}
\lean{Graph/LocalRewiring}{17}{19}{Linear.Focussed.movePort}
Joining a port equals contracting it and then lifting the retained port back to
original names. If it is the focus target, both sides yield the focus source;
otherwise contraction followed by lifting fixes it. Moreover $M_w(v)=v$ when
$v\ne b$ and $v\ne J(w)$: both tests in the displayed definitions are false.
\lean{Graph/LocalRewiring}{21}{31}{Linear.Focussed.joinPort_eq_lift_contract}
\lean{Graph/LocalRewiring}{34}{39}{Linear.Focussed.movePort_unchanged}

\begin{definition}[Local rewiring]\label{term:local-rewiring}
For states $s,t$ and a selected pair $(m,w)$, a local rewiring comprises
bijections $V$ and $E$ of their internal vertex and edge carriers. Extend $V$
to ports by fixing the ordered boundary. The bijections preserve internal
vertex labels and identify the distinguished focus edges. The target focus mode
is $m$, and every non-focus edge retains its entire label, including box bodies.
The target incidence lists obey
\begin{align*}
 t.\mathsf{source}(E(f))&=[V(J(w))],\\
 t.\mathsf{source}(E(e))&=\mathsf{map}(V\circ M_w,s.\mathsf{source}(e))
       &&(e\ne f),\\
 t.\mathsf{target}(E(e))&=\mathsf{map}(V,s.\mathsf{target}(e))
       &&\text{for every }e.
\end{align*}
Here $f$ is the source state's focus edge. The focus label follows from the
target state's focus certificate and its mode equation. Well-formedness of
both states is part of their types; these equations do not postulate that an
arbitrary raw graph patch is well formed.
\end{definition}
\lean{Graph/LocalRewiring}{46}{60}{Linear.LocalRewiring}
The private port-composition lemma says that extending a composite vertex
bijection to ports equals composing its two extensions. Check the internal,
input, and output constructors separately; each equality is immediate.
\lean{Graph/LocalRewiring}{64}{68}{Linear.LocalRewiring.port_trans}

\subsection{Canonical construction and uniqueness}
Fix a typed erased coordinate $q=(m,\mathsf{contract}(w))$. Reinsertion at $q$
introduces a singleton vertex and edge. The restoration bijections send those
fresh elements to the old focus target and old focus edge, and retain all
other names.
\begin{lemma}[Restoration implements port movement]
For every original port $v$, restoring the names of the split contracted port
at $q$ yields $M_w(v)$.
\end{lemma}
\begin{proof}
Injectivity of lifting shows that the contracted ports of $v$ and $w$ agree
exactly when $J(v)=J(w)$. In that case insertion splits to the fresh vertex,
which restoration sends to the old focus target. Otherwise insertion retains
the contracted port; restoring that host port yields its lift, namely $J(v)$.
These are precisely the two defining cases of $M_w$.
\end{proof}
\lean{Graph/LocalRewiring}{75}{96}{Linear.LocalRewiring.restore_move}

\begin{theorem}[Characterization of a rewired target]
Every target satisfying the local equations is focus-preservingly isomorphic
to the canonical state realized from $q$. Conversely the canonical state
satisfies the local equations. Replacing a target by a focus-preservingly
isomorphic state preserves local rewiring.
\end{theorem}
\begin{proof}
For the first direction compose restoration with the given carrier bijections.
On a retained vertex its label equation is the supplied one; on the fresh
vertex use that the old focus target has term type. Retained edge labels use
the non-focus label equation; the fresh edge uses the target focus's label
and the selected mode. For retained sources apply the restoration lemma
pointwise. For the focus source use joining equals lifted contraction.
For retained targets, contraction followed by lifting is the identity since
no non-focus edge targets the deleted focus target. The fresh target restores
to that target. These equations provide every field of the focussed isomorphism.

For the converse use the inverse restoration bijections. The existing
reinsertion isomorphism supplies internal-label and target equations, which
are independent of the new coordinate and mode. Non-focus labels are unchanged
by construction. Apply injectivity of the forward restoration port map to the
source equations: the retained-source equation becomes the restoration lemma,
and the focus-source equation becomes joining equals lifted contraction.
The focus edge and mode equations follow from the definitions.

Finally compose carrier bijections with a target isomorphism. Its incidence
preservation, list-map composition, and the port-composition lemma prove all
new source and target equations. Labels, mode, and focus identity compose
in the same order.
\end{proof}
\lean{Graph/LocalRewiring}{100}{160}{Linear.LocalRewiring.iso}
\lean{Graph/LocalRewiring}{163}{207}{Linear.LocalRewiring.canonical}
\lean{Graph/LocalRewiring}{210}{231}{Linear.LocalRewiring.transport}

\subsection{State-level equivalence and regression examples}
\begin{definition}[Incidence search]\label{term:incidence-search}
An incidence-search step $s\to t$ consists of a selected pair $(m,w)$, an
original-graph local interaction selecting it, and a local rewiring from $s$
to $t$ at that pair. This definition contains no coordinate-search premise.
\end{definition}
\lean{Graph/LocalRewiring}{237}{238}{Linear.State.IncidenceSearch}
\begin{theorem}[Exact correspondence on actual states]
Erasure-based state search is equivalent to incidence search. Recursive state
search is equivalent to an incidence step followed by recursive focussed
isomorphism. Incidence search is deterministic up to literal-label focussed
isomorphism.
\end{theorem}
\begin{proof}
From erasure-based search, recover an original local interaction, construct the
canonical rewiring, and transport it through the supplied target isomorphism.
Conversely a local interaction gives a coordinate step. Type preservation of
coordinate search certifies its target as a typed position; the rewiring
characterization supplies the required target isomorphism. The recursive
statement follows by substituting this equivalence into the existing
characterization of recursive search as literal search plus recursive endpoint
renaming. Determinism follows by converting two incidence steps to state
search and applying state-search determinism.
\end{proof}
\lean{Graph/LocalRewiring}{241}{250}{Linear.State.search_iff_incidence}
\lean{Graph/LocalRewiring}{252}{256}{Linear.State.deepSearch_iff_incidence}
\lean{Graph/LocalRewiring}{259}{263}{Linear.State.incidenceSearch_deterministic}

The examples recover the actual nullary value edge in the initial value state.
Its source is the focus target and its argument and thunk lists are empty;
the source incidence follows by splitting the unique input port. Direct entry
therefore selects backtracking at the old target. This exercises the joining
case $J(b)=a$, not merely an unaffected port. The known terminating value step
converts to an incidence step whose target is final.
\lean{Graph/LocalSearchExamples}{9}{21}{Linear.LocalSearchExamples.valueOperation}
\lean{Graph/LocalSearchExamples}{25}{28}{Linear.LocalSearchExamples.value_interaction}
\lean{Graph/LocalSearchExamples}{31}{34}{Linear.LocalSearchExamples.value_terminates}

A second example inserts an initial search focus before the variable in the
closed variable/value chain. Its variable source list is checked directly to
be the singleton focus target, so the local variable rule selects rewrite
mode at the old source. Type preservation and canonical rewiring construct
an actual incidence successor in rewrite mode. This proves only the trigger,
not a substitution transition.
\lean{Graph/LocalSearchExamples}{37}{39}{Linear.LocalSearchExamples.initialChain}
\lean{Graph/LocalSearchExamples}{42}{48}{Linear.LocalSearchExamples.variable_interaction}
\lean{Graph/LocalSearchExamples}{52}{60}{Linear.LocalSearchExamples.variable_successor}
